%----------------------------------------------------------------------------------------
% FORMULES OMVORMEN. Alle opgaven zijn eigen werk.
%
% Opbouw: Probleem, Start, leerstof, en op een nieuwe pagina de Opdrachten en Test jezelf
% (2de graad: geen Koppel). Labels van dit hoofdstuk beginnen met h3-.
% Pakket 4 en 5: ook formules waarin de letter twee keer staat, en meer formules uit de fysica.
%----------------------------------------------------------------------------------------

\mychapter{3}{Formules omvormen}

\begin{probleem}
Je plant een reis naar New York en bekijkt het weerbericht. Daar meten ze de temperatuur in graden Fahrenheit. Voor de komende vijf dagen voorspellen ze $41$, $50$, $59$, $68$ en $77\,^\circ$F. Je weet dat een temperatuur van $C$ graden Celsius gelijk is aan
\[ F = 1{,}8C + 32 \]
graden Fahrenheit. Hoe warm wordt het in graden Celsius? Moet je daarvoor vijf keer een vergelijking oplossen?
\end{probleem}

\lkr{Leerplan \pakketnaam: LPD \alleenvoor{1,2,3}{16}\alleenvoor{4,5}{20} (formules omvormen). Voorkennis: hoofdstuk 1, en de vierkantswortel uit de 1ste graad (in sectie 2). Laat leerlingen het probleem eerst getal per getal oplossen: na twee of drie vergelijkingen zien ze dat ze telkens hetzelfde doen. Dat is het idee van omvormen. \alleenvoor{4,5}{In dit pakket komen ook formules aan bod waarin de letter twee keer staat, en een paar formules extra uit de fysica.}}

Formules kom je overal tegen: in de fysica, in de economie, in de meetkunde, op een factuur. Een formule zegt hoe je één grootheid berekent uit andere. Maar soms ken je net die ene grootheid, en zoek je een van de andere. Dan los je een vergelijking op. In dit hoofdstuk leer je dat in één keer te doen, met letters in plaats van getallen: een formule omvormen. Je gebruikt daarvoor precies de rekenregels van hoofdstuk 1.

\begin{oefeningen}[Start]
\oef{h3-start-1}

\oef{h3-verken}
\end{oefeningen}

%========================================================================================
\section{Een formule omvormen}

\subsection{Van getallen naar letters}

In de Verken zocht je de breedte $b$ van een rechthoek met oppervlakte $A = l \cdot b$. Met $A = 24$ en $l = 6$ los je $24 = 6b$ op, en vind je $b = \tfrac{24}{6} = 4$. Met andere getallen deel je opnieuw $A$ door $l$. Dat kun je één keer doen, met de letters zelf.

\voorbeeld{}{Vorm de formule $A = l \cdot b$ om naar $b$.\label{vb:h3-verken}}{%
\uitleg{We doen met de letters precies wat we met de getallen deden. Links staat de berekening met getallen, rechts dezelfde stap met letters.}
{\renewcommand{\arraystretch}{1.6}%
\begin{tabular}{@{}>{$}l<{$}@{\hspace{1.4cm}}>{$}l<{$}@{\hspace{1cm}}>{\small\itshape}l@{}}
24 = 6 \cdot b & A = l \cdot b & De opgave. \\
\dfrac{24}{6} = b & \dfrac{A}{l} = b & Deel beide leden door het getal of de letter voor $b$. \\
b = 4 & \formule{b = \dfrac{A}{l}} & De formule naar $b$.
\end{tabular}\par}
Controle met de getallen: $\tfrac{24}{6} = 4$. Met de nieuwe formule vind je de breedte meteen, voor elke rechthoek.}

Zo'n nieuwe formule heet de formule \textit{omgevormd naar} $b$. Ze bevat dezelfde informatie als $A = l \cdot b$, maar ze is handiger als je $b$ zoekt.

\begin{theorie}
Een formule \textbf{\textcolor{accent}{omvormen}} naar een letter is ze zo herschrijven dat die letter alleen in één lid staat, en niet in het andere lid.

\vspace{1mm}
Je gebruikt de rekenregels van vergelijkingen. De andere letters behandel je als getallen. Delen door een letter mag alleen als die letter niet 0 is.
\end{theorie}

In formules uit de praktijk zijn de grootheden meestal positief: een lengte, een tijd, een massa. Delen door $l$ is dan geen probleem.

\voorbeeld{}{De afgelegde weg $s$ bij een snelheid $v$ gedurende een tijd $t$ is $s = v \cdot t$. Vorm om naar $t$.\label{vb:h3-t}}{%
\uitleg{De letter $t$ staat vermenigvuldigd met $v$. We delen beide leden door $v$.}
\begin{stappen}
\dfrac{s}{v} = t & Deel beide leden door $v$, met $v \neq 0$. \\[2ex]
t = \dfrac{s}{v} & Zet $t$ links: dat leest makkelijker.
\end{stappen}
Controle met getallen: wie 120~km aflegt met 80~km/h, rijdt $\tfrac{120}{80} = 1{,}5$~uur.}

Een formule mag je altijd omdraaien: $\tfrac sv = t$ en $t = \tfrac sv$ zeggen hetzelfde. Meestal zet je de gezochte letter links.

\voorbeeld{}{Vorm $v = \dfrac{s}{t}$ om naar $s$.\label{vb:h3-s}}{%
\uitleg{Nu staat $s$ in een teller, gedeeld door $t$. Het omgekeerde van delen door $t$ is vermenigvuldigen met $t$.}
\begin{stappen}
v \cdot t = s & Vermenigvuldig beide leden met $t$. \\
s = v \cdot t & Zet $s$ links.
\end{stappen}}

\begin{werkwijze}
Kijk welke bewerkingen er met de gezochte letter gebeuren. Maak ze één voor één ongedaan met de rekenregels van vergelijkingen, in beide leden: een optelling met een aftrekking, een vermenigvuldiging met een deling, en omgekeerd. Controleer met eenvoudige getallen.
\end{werkwijze}

\nujij{Vorm $U = R \cdot I$ om naar $I$.}{Deel door $R$: $I = \tfrac{U}{R}$.}

\begin{oefeningen}
\oef{h3-eenvoudig}
\end{oefeningen}

\subsection{Twee stappen}

Wat als er meer gebeurt met de gezochte letter? Een fiets huren voor $u$ uur kost $P = 3u + 5$ euro. Daar wordt $u$ eerst met 3 vermenigvuldigd, en daarna wordt er 5 bij opgeteld. Om $u$ terug te vinden, maak je die stappen ongedaan in de omgekeerde volgorde: eerst 5 aftrekken, dan delen door 3. Zoals bij aankleden: wat je laatst aantrekt, trek je eerst uit. Zo vind je $u = \tfrac{P - 5}{3}$.

\voorbeeld{}{Terug naar het probleem. Je plant een reis naar New York. Daar meten ze de temperatuur in graden Fahrenheit, en voor de komende vijf dagen voorspellen ze $41$, $50$, $59$, $68$ en $77\,^\circ$F. Een temperatuur van $C$ graden Celsius is $F = 1{,}8C + 32$ graden Fahrenheit. Hoe warm wordt het in graden Celsius? Vorm daarvoor de formule om naar $C$, in plaats van vijf keer een vergelijking op te lossen.\label{vb:h3-terug}}{%
\uitleg{We maken de laatste bewerking eerst ongedaan: de optelling met 32. Daarna de vermenigvuldiging met $1{,}8$.}
\naast{\begin{stappenvrij}{3.4cm}{4.4cm}
F - 32 = 1{,}8C & Trek 32 af van beide leden. \\[1ex]
\dfrac{F - 32}{1{,}8} = C & Deel beide leden door $1{,}8$. \\[2ex]
C = \dfrac{F - 32}{1{,}8} & De formule naar $C$.
\end{stappenvrij}

\vspace{2mm}
{\small\renewcommand{\arraystretch}{1.3}%
\begin{tabular}{|l|c|c|c|c|c|}
\hline
$F$ (in $^\circ$F) & 41 & 50 & 59 & 68 & 77 \\ \hline
$C$ (in $^\circ$C) & 5 & 10 & 15 & 20 & 25 \\ \hline
\end{tabular}\par}}{\begin{tikzpicture}[x=.09cm,y=.028cm,font=\scriptsize]
  \draw[black!14,very thin,xstep=10,ystep=20] (0,0) grid (30,100);
  \draw[->,grijs] (0,0) -- (33,0) node[right,text=black] {$C$};
  \draw[->,grijs] (0,0) -- (0,108) node[above,text=black] {$F$};
  \foreach \t in {10,20,30}{\node[below,inner sep=1.5pt] at (\t,0) {$\t$};}
  \foreach \t in {20,40,60,80,100}{\node[left,inner sep=1.5pt] at (0,\t) {$\t$};}
  \draw[accent,thick] (0,32) -- (30,86);
  \foreach \c in {5,10,15,20,25}{\fill[accent] (\c,{1.8*\c+32}) circle (1.6pt);}
\end{tikzpicture}}
Bij $41\,^\circ$F is het dus $\tfrac{41 - 32}{1{,}8} = \tfrac{9}{1{,}8} = 5\,^\circ$C. Je hoeft niet vijf keer een vergelijking op te lossen: één keer omvormen volstaat, daarna vul je alleen nog in. Het wordt een lentereis van $5$ tot $25\,^\circ$C.}

De figuur toont het verband: elke $^\circ$C erbij geeft $1{,}8\,^\circ$F erbij, en bij $0\,^\circ$C is het al $32\,^\circ$F. Het verband werkt in beide richtingen: zo lees je ook in de figuur af dat $68\,^\circ$F gelijk is aan $20\,^\circ$C.

Soms kun je kiezen in welke volgorde je werkt. Beide manieren zijn juist, als je maar in beide leden hetzelfde doet.

\voorbeeld{}{De omtrek van een rechthoek is $P = 2(l + b)$. Vorm om naar $b$.\label{vb:h3-omtrek}}{%
\uitleg{Hier kun je eerst de haakjes uitwerken, of eerst door 2 delen. We doen het op de tweede manier, en vergelijken met de eerste.}
\begin{stappen}
\dfrac{P}{2} = l + b & Deel beide leden door 2. \\[2ex]
b = \dfrac{P}{2} - l & Trek $l$ af van beide leden. \\[2ex]
P = 2l + 2b \ \Rightarrow \ b = \dfrac{P - 2l}{2} & De andere manier: werk de haakjes uit, trek $2l$ af en deel door 2.
\end{stappen}
Beide formules zijn gelijk, want $\tfrac{P - 2l}{2} = \tfrac P2 - \tfrac{2l}{2} = \tfrac P2 - l$. Controle met $P = 20$ en $l = 6$: $b = 10 - 6 = 4$, en $2(6 + 4) = 20$.}

Een omgevormde formule kun je altijd controleren met eenvoudige getallen. Kies waarden die bij de oorspronkelijke formule horen, en kijk of de nieuwe formule hetzelfde geeft.

\nujij{Vorm $y = 3x - 6$ om naar $x$.}{$y + 6 = 3x$, dus $x = \tfrac{y + 6}{3}$.}

\begin{oefeningen}
\oef{h3-tweestappen}

\oef{h3-controle}
\end{oefeningen}

%========================================================================================
\section{Breuken en kwadraten}

Niet elke formule is zo eenvoudig als $s = v \cdot t$. In de meetkunde en de fysica staan er vaak breuken in, of kwadraten. Het idee blijft hetzelfde: maak de bewerkingen op de gezochte letter één voor één ongedaan.

\subsection{Een breuk in de formule}

\naast{De oppervlakte van een trapezium met evenwijdige zijden $B$ en $b$ en hoogte $h$ is
\[ A = \dfrac{(B + b) \cdot h}{2} . \]
Hoe hoog is een trapezium met een gegeven oppervlakte? Werk eerst de noemer weg, net zoals bij vergelijkingen met breuken.}{\begin{tikzpicture}[scale=.42,font=\scriptsize]
  \draw[accent,thick,fill=accent!10] (0,0) -- (7,0) -- (5.2,3) -- (1.4,3) -- cycle;
  \draw[grijs,dashed] (1.4,3) -- (1.4,0);
  \node[below] at (3.5,0) {$B$};
  \node[above] at (3.3,3) {$b$};
  \node[left] at (1.4,1.5) {$h$};
\end{tikzpicture}}

\voorbeeld{}{Vorm $A = \dfrac{(B + b) \cdot h}{2}$ om naar $h$.\label{vb:h3-trapezium}}{%
\uitleg{De letter $h$ wordt vermenigvuldigd met $(B + b)$, en daarna gedeeld door 2. We maken dat ongedaan in de omgekeerde volgorde.}
\begin{stappen}
2A = (B + b) \cdot h & Vermenigvuldig beide leden met 2. \\
h = \dfrac{2A}{B + b} & Deel beide leden door $B + b$. Die som staat in zijn geheel in de noemer.
\end{stappen}
Controle: een trapezium met $B = 7$, $b = 3$ en $h = 4$ heeft $A = \tfrac{10 \cdot 4}{2} = 20$, en $\tfrac{2 \cdot 20}{7 + 3} = 4$.}

\voorbeeld{}{Vorm dezelfde formule om naar $b$.\label{vb:h3-trapezium2}}{%
\uitleg{Nu zit $b$ in de som $B + b$. We maken eerst de deling door 2 ongedaan, dan de vermenigvuldiging met $h$, en dan de optelling met $B$.}
\begin{stappen}
2A = (B + b) \cdot h & Vermenigvuldig beide leden met 2. \\[1ex]
\dfrac{2A}{h} = B + b & Deel beide leden door $h$. \\[2ex]
b = \dfrac{2A}{h} - B & Trek $B$ af.
\end{stappen}
Controle met de getallen hierboven: $\tfrac{2 \cdot 20}{4} - 7 = 10 - 7 = 3$.}

\nujij{De oppervlakte van een driehoek is $A = \dfrac{b \cdot h}{2}$. Vorm om naar $h$.}{$2A = b \cdot h$, dus $h = \tfrac{2A}{b}$.}

\begin{oefeningen}
\oef{h3-breuk}
\end{oefeningen}

\subsection{De letter in de noemer}

De dichtheid van een stof is de massa per volume: $\rho = \tfrac{m}{V}$, met de Griekse letter $\rho$ (\textit{rho}). Hoeveel ruimte neemt 1~kg ijzer in? Dan zoek je $V$, en die staat in de noemer.

\voorbeeld{}{Vorm $\rho = \dfrac{m}{V}$ om naar $V$.\label{vb:h3-noemer}}{%
\uitleg{Staat de gezochte letter in de noemer, haal haar daar dan eerst uit: vermenigvuldig beide leden met de noemer.}
\begin{stappen}
\rho \cdot V = m & Vermenigvuldig beide leden met $V$. Nu staat $V$ niet meer in een noemer. \\
V = \dfrac{m}{\rho} & Deel beide leden door $\rho$.
\end{stappen}
IJzer heeft een dichtheid van $7{,}9$~g/cm$^3$. Dus 1~kg, of 1000~g, ijzer neemt $\tfrac{1000}{7{,}9} \approx 127$~cm$^3$ in.}

Een veelgemaakte fout is $V = \rho \cdot m$. Controleer met getallen: $\tfrac{12}{3} = 4$, dus moet $3 = \tfrac{12}{4}$, en niet $3 = 4 \cdot 12$.

\begin{oefeningen}
\oef{h3-noemer}
\end{oefeningen}

\subsection{Een kwadraat}

\naast{De oppervlakte van een cirkel met straal $r$ is $A = \pi r^2$. Hoe groot is de straal van een rond tapijt van $3~\text{m}^2$? Nu staat $r$ in het kwadraat. Het omgekeerde van kwadrateren ken je uit de 1ste graad: de vierkantswortel. Zo is $\sqrt{25} = 5$, want $5^2 = 25$. Ook $(-5)^2 = 25$, maar een straal is positief: de positieve wortel is de juiste.}{\begin{tikzpicture}[scale=.6,font=\scriptsize]
  \fill[accent!10] (0,0) circle (1.6);
  \draw[accent,thick] (0,0) circle (1.6);
  \draw[grijs] (0,0) -- (1.6,0) node[midway,above] {$r$};
  \fill (0,0) circle (1.5pt);
\end{tikzpicture}}

\voorbeeld{}{Vorm $A = \pi r^2$ om naar $r$, met $r > 0$.\label{vb:h3-kwadraat}}{%
\uitleg{Eerst halen we $\pi$ weg, daarna trekken we de vierkantswortel.}
\begin{stappen}
\dfrac{A}{\pi} = r^2 & Deel beide leden door $\pi$. \\[2ex]
r = \sqrt{\dfrac{A}{\pi}} & Neem de positieve vierkantswortel, want $r > 0$.
\end{stappen}
Voor het tapijt is $r = \sqrt{\tfrac{3}{\pi}} \approx 0{,}98$~m.}

Let op: de wortel staat over de hele breuk. $\tfrac{\sqrt{A}}{\pi}$ is iets anders. Controle: een cirkel met $r = 2$ heeft $A = 4\pi$, en $\sqrt{\tfrac{4\pi}{\pi}} = \sqrt4 = 2$.

\begin{oefeningen}
\oef{h3-kwadraat}
\end{oefeningen}

\begin{alleenvoor}{4,5}
\subsection{De letter staat twee keer}

Zet je geld op een spaarrekening met een enkelvoudige intrest van $i$ per jaar, dan heb je na $n$ jaar een eindkapitaal $K = k + k \cdot i \cdot n$, met $k$ het beginkapitaal. Hoeveel moet je vandaag sparen om later een bepaald bedrag te hebben? Dan zoek je $k$, en die letter staat twee keer in de formule. Delen door $k$ helpt niet, want dan blijft er $k$ staan. Het idee: breng alle termen met $k$ samen, en zet $k$ buiten haakjes.

\voorbeeld{}{Vorm $K = k + k \cdot i \cdot n$ om naar $k$.\label{vb:h3-twee}}{%
\uitleg{Beide termen in het rechterlid bevatten $k$. We zetten $k$ buiten haakjes: dan staat $k$ nog maar één keer.}
\begin{stappen}
K = k\,(1 + i \cdot n) & Zet $k$ buiten haakjes: $k \cdot 1 + k \cdot in$. \\
k = \dfrac{K}{1 + i \cdot n} & Deel door $1 + in$, en dat is positief.
\end{stappen}
Wie over 3 jaar 2240 euro wil hebben, bij een intrest van 4\,\% per jaar, moet nu $\tfrac{2240}{1 + 0{,}04 \cdot 3} = \tfrac{2240}{1{,}12} = 2000$ euro sparen.}

\begin{werkwijze}
Staat de gezochte letter in meer dan één term, breng die termen dan naar één lid, en de andere termen naar het andere lid. Zet de letter buiten haakjes, en deel door wat tussen de haakjes staat.
\end{werkwijze}

\begin{oefeningen}
\oef{h3-twee}
\end{oefeningen}
\end{alleenvoor}

%========================================================================================
\section{Formules gebruiken}

Wanneer vorm je een formule om, en wanneer vul je gewoon de getallen in en los je de vergelijking op? Beide geven hetzelfde antwoord. Moet je maar één keer rekenen, dan is invullen vaak het snelst. Moet je dezelfde berekening vaak doen, bijvoorbeeld voor een hele tabel van waarden, of wil je de formule in een rekenblad zetten, dan loont omvormen.

\voorbeeld{}{Een bakker verkoopt taarten voor $p$ euro per stuk. De ingrediënten kosten $k$ euro per taart, en elke maand betaalt ze $V$ euro vaste kosten (huur, elektriciteit). Haar winst bij $q$ verkochte taarten is
\[ W = (p - k) \cdot q - V . \]
Vorm om naar $q$. Hoeveel taarten moet ze per maand verkopen om quitte te spelen, met $p = 18$, $k = 6$ en $V = 2400$? En voor 1200 euro winst?\label{vb:h3-winst}}{%
\uitleg{Op $q$ gebeuren twee bewerkingen: vermenigvuldigen met $(p - k)$, en dan $V$ aftrekken. Die maken we ongedaan in de omgekeerde volgorde. Quitte spelen betekent: winst 0.}
\begin{stappen}
W + V = (p - k) \cdot q & Tel $V$ op bij beide leden. \\
q = \dfrac{W + V}{p - k} & Deel door $p - k$, de winst per taart zonder de vaste kosten. \\[2ex]
q = \dfrac{0 + 2400}{18 - 6} = 200 & Quitte spelen: vul $W = 0$ in. \\[2ex]
q = \dfrac{1200 + 2400}{12} = 300 & Voor 1200 euro winst.
\end{stappen}
\textit{Ze moet 200 taarten verkopen om quitte te spelen, en 300 voor 1200 euro winst.} Elke taart boven de 200 brengt 12 euro winst op.}

Let bij het invullen op de eenheden. In een formule uit de fysica moeten ze bij elkaar passen: een snelheid in km/h hoort bij een afstand in km en een tijd in uur.

\begin{alleenvoor}{4,5}
\voorbeeld{}{Een voorwerp met massa $m$ (in kg) en snelheid $v$ (in m/s) heeft een bewegingsenergie $E = \tfrac12 m v^2$ (in joule). Hoe snel rijdt een auto van 1200~kg met een bewegingsenergie van 240\,000~J?\label{vb:h3-energie}}{%
\uitleg{We vormen om naar $v$, met $v > 0$. Daarna vullen we in.}
\begin{stappen}
2E = m v^2 & Vermenigvuldig beide leden met 2. \\[1ex]
\dfrac{2E}{m} = v^2 \ \Rightarrow \ v = \sqrt{\dfrac{2E}{m}} & Deel door $m$ en neem de positieve wortel. \\[2ex]
v = \sqrt{\dfrac{480\,000}{1200}} = \sqrt{400} = 20 & Vul in.
\end{stappen}
\textit{De auto rijdt 20~m/s,} of $20 \cdot 3{,}6 = 72$~km/h.}
\end{alleenvoor}

\begin{werkwijze}
Zoek de formule die de grootheden verbindt. Vorm ze om naar de gezochte grootheid. Kijk of de eenheden bij elkaar passen, vul de gegevens in en reken uit. Controleer of de uitkomst kan, en antwoord met een zin en een eenheid.
\end{werkwijze}

\begin{oefeningen}
\oef{h3-gebruik}
\end{oefeningen}

%========================================================================================
\opdrachtenkop

\begin{oefeningen}
\oef{h3-o1}

\oef{h3-o2}

\oef{h3-o3}

\oef{h3-o4}

\oef{h3-o5}

\oef{h3-o6}

\oef{h3-o7}

\oef{h3-o8}

\oef{h3-o9}

\oef{h3-o10}

\oef{h3-o11}

\oef{h3-o12}

\oef{h3-o13}

\oef{h3-o14}

\oef{h3-o15}

\oef{h3-o16}

\oef{h3-o17}

\oef{h3-o18}

\oef{h3-o19}

\oef{h3-o20}

\oef{h3-o21}

\oef{h3-o22}

\oef{h3-o23}

\oef{h3-o24}
\end{oefeningen}

%----------------------------------------------------------------------------------------
\needspace{14\baselineskip}
\opdrachtsectie{Test jezelf}

\begin{oefeningen}
\oef{h3-t1}

\oef{h3-t2}

\oef{h3-t3}

\oef{h3-t4}

\oef{h3-t5}

\oef{h3-t6}

\oef{h3-t7}

\oef{h3-t8}
\end{oefeningen}

\herhaalschema{%
  \fv{oef:h3-t1}{vb:h3-t}{oef:h3-eenvoudig}%
  \fv{oef:h3-t2}{vb:h3-omtrek}{oef:h3-tweestappen}%
  \fv{oef:h3-t3}{vb:h3-trapezium}{oef:h3-breuk}%
  \fv{oef:h3-t4}{vb:h3-noemer}{oef:h3-noemer}%
  \fv{oef:h3-t5}{vb:h3-kwadraat}{oef:h3-kwadraat}%
  \fv{oef:h3-t6}{vb:h3-terug}{oef:h3-o9}%
  \fv{oef:h3-t7}{vb:h3-winst}{oef:h3-gebruik}%
  \alleenvoor{4,5}{\fv{oef:h3-t8}{vb:h3-twee}{oef:h3-twee}}}

\ifzelfstudie
\section*{Maak je eigen samenvatting}
Sluit dit boek en schrijf in je schrift: wat omvormen betekent, waarom je de bewerkingen in de omgekeerde volgorde ongedaan maakt, wat je doet met een breuk, met een letter in de noemer en met een kwadraat, \alleenvoor{4,5}{wat je doet als de letter twee keer staat, }en hoe je een omgevormde formule controleert met getallen. Vergelijk daarna met de kaders.
\fi
